\section{Threats to Validity}
\label{sec:threats}

Each item says how it is known: \emph{measured} (a number from the data),
\emph{recorded} (a fact of the setup) or \emph{not established} (a likely
cause that no data here isolates).

\subsection{Internal validity}

\begin{itemize}[leftmargin=1.5em]
\item \emph{One P-core is degraded.} Measured: on the prepared machine it
gives the right checksum and the time of the other cores, and it took 9 of
the 30 native one-thread filter runs with a median 1.0003 times the
configuration's. Not established: that it never slows down within a run;
ten pinned runs in one session do not rule out a core that is slow now and
then. At $N=8$ every run of every environment uses it.
\item \emph{Thread placement is left to the kernel}, because a browser
cannot pin threads. Measured: the shared-atomic cells at $k=1$ land on
several time levels from run to run, which makes them the widest cells
(Section~\ref{sec:levels}). Not established: what sets the level of a run.
\item \emph{The layers below the crate differ} (Section~\ref{sec:builds}):
WASM-only flags, a rebuilt standard library, other locks inside Rayon, idle
workers that spin instead of yielding, and the engines' compilers. They are
the environment under study; ``the same code'' is claimed for the crate
only.
\item \emph{The compiler tier.} Measured: in every cell of the tier
campaign the browser as shipped equals its optimizing-only arm after the
warm-up, so the campaigns measure optimized code. Recorded: the listings of
V8's code come from diagnostic runs with other flags.
\item \emph{Slow early iterations in one Firefox cell.} Measured: one or
two of the first three measured iterations of the local histogram at $k=1$,
$N=8$ are about 10\% slower, after 71 warm-up iterations, in 27 of 30 runs;
the run value moves by about 0.1\%. This cell is the control of one primary
contrast. Not established: what they are; they come with the optimizing
compiler's code.
\item \emph{The fork-join calibration has 20 samples and a coarse clock.}
Recorded: with a 20~\textmu s clock Firefox's warm fork-join reads 0 at
$N \leq 4$.
\item \emph{Browser threads share the measured cores.} Recorded: Firefox
reports a hardware concurrency of 32 and Chrome of 8, so Firefox sizes its
helper threads for the whole machine while confined to eight cores. Not
established: how much CPU, cache or memory bandwidth they take during a
timed iteration. With the start delay no browser shows a knee at eight
threads.
\item \emph{A coding assistant was alive during the campaigns}, confined to
the E-cores under a no-command rule. Measured: load and foreign CPU time
equal those of a campaign without it. Not established: that nothing ran on
the other cores that shares the cache.
\item \emph{Frequency is not known during a run.} Recorded: turbo off and
the governor fixed. Measured: the kernel's frequency sample read 1154 to
3201~MHz on cores that had just run at full speed, so it cannot tell a slow
core from one read while going idle; it is recorded and not used.
\item \emph{Memory growth.} Measured: the native filter's memory grows by
exactly 65\,536~KiB in the warm-up and by nothing in the measured phase at
one and two threads, and by a few KiB (stack pages) with more.
\item \emph{The native runs have no start delay} and the browsers wait 3~s
idle. It affects $T_1$ and $T_N$ of an environment alike.
\item \emph{The counters of the loop} come from two campaigns run one after
the other, so the pass count is confounded with the block.
\end{itemize}

\subsection{Construct validity}

\begin{itemize}[leftmargin=1.5em]
\item \emph{RQ2 is a ratio of ratios.} Measured: the environment gap of the
local control moves it as much as the atomics do
(Section~\ref{sec:decomposition}); it is reported with both factors and
with the ratio of the excesses.
\item \emph{The control is not neutral in $k$.} At small $k$ the local
histogram is itself a dependent chain. The contrast compares within one
$k$, never across.
\item \emph{Speedup is relative to a pool of one thread.} Measured: that
pool is within $-7$\% to $+5$\% of sequential code in every environment
(Section~\ref{sec:sequential}); the time over native and the absolute
speedup are reported next to it.
\item \emph{The efficiency dip at $N=6$ and 7 is not explained.} Measured:
it appears in every environment, and fixed grains do not remove it
(Section~\ref{sec:dip}).
\item \emph{The atomic code differs between engines}, which is the finding
of RQ2 and also a limit of its name: what is compared is ``shared atomic
counters as each compiler emits them'', not one instruction in three places.
\item \emph{The cutoff sizes are lower bounds for a single call}
(Section~\ref{sec:cutoff}).
\end{itemize}

\subsection{Statistical conclusion validity}

\begin{itemize}[leftmargin=1.5em]
\item \emph{Precision.} With 30 runs the primary 99.5\% intervals are
within $\pm$2\% for RQ1 and $\pm$7\% for RQ2 (the widest, Chrome at $k=1$).
The two non-rejections are bounded by their intervals, about 1\% and 6\%,
and are not equalities.
\item \emph{Percentile bootstrap intervals over 30 runs} run slightly
narrow at the 99.5\% level; the plan fixed them before the data.
\item \emph{The points of one speedup curve share their $T_1$}, so an error
in $T_1$ shifts the whole curve.
\item \emph{The smallest rejected effect is 3\%} (Firefox at $k=4$). On the
earlier build the Firefox values at $k=4$ and 256 moved between two builds
by more than that, so this rejection is not read as an effect of the
environment; the reading of Firefox rests on the two factors.
\item \emph{This is a replication whose first results were known.} The
rule for adopting the build read only counts of runs, and the analysis
plan did not change.
\item \emph{Exploratory results} use uncorrected 95\% intervals and decide
nothing. The robustness check was added after the first primary results
existed.
\end{itemize}

\subsection{External validity}

\begin{itemize}[leftmargin=1.5em]
\item \emph{One machine, one compiler version, one version of each
browser.} The RQ2 finding is about the code V8 155 emits for an atomic add;
another version may emit another sequence, and the same engine with its
other compiler already does.
\item \emph{P-cores only, turbo off.} The hybrid campaign shows what
E-cores add; nothing here describes turbo.
\item \emph{Every primary cell lives in memory, not in cache}: 128~MiB for
the filter's two images and 64~MiB for the histograms, against a 36~MB
last-level cache. The explanation of RQ1 is about this regime. The in-cache
regime appears only in the cutoff campaigns.
\item \emph{Inputs are uniform random values} with no spatial structure,
one seed, one size.
\item \emph{Two workloads}, both data-parallel loops over an image. Nothing
here covers locks, channels, allocation-heavy code (beyond one sort) or
irregular parallelism.
\item \emph{Headless browsers} on an idle machine; a pilot bounds the
difference to headful at 0.3\%.
\end{itemize}
